\documentclass[11pt,a4paper]{article}
\usepackage[utf8]{inputenc}
\usepackage[T1]{fontenc}
\usepackage{amsmath,amssymb}
\usepackage{booktabs}
\usepackage{geometry}
\usepackage{enumitem}
\usepackage{xcolor}
\usepackage{hyperref}
\geometry{margin=2.3cm}
\definecolor{accent}{RGB}{38,70,83}
\definecolor{ours}{RGB}{42,157,143}
\hypersetup{colorlinks=true,linkcolor=accent,urlcolor=accent}
\newcommand{\code}[1]{\texttt{\small #1}}

\title{\vspace{-1.5em}\textbf{Document 3 --- Model Architectures}\\
\large Lightweight Students, the U-Net Variant, and the Mamba Variant}
\author{}
\date{}

\begin{document}
\maketitle
\tableofcontents
\vspace{1em}\hrule\vspace{1em}

All models map the same input/output: $[B,4,161,361]\!\to\![B,1,161,361]$, a 4-AP heatmap stack
to a single location heatmap, finished with a Sigmoid. They differ only in \emph{how} they
compress and reconstruct the spatial information. Table~\ref{tab:summary} summarizes the family;
the sections that follow give layer-level detail.

\begin{table}[h]
\centering
\caption{Model family. Compression is relative to the 16.5\,M-parameter DLoc teacher.}
\label{tab:summary}
\begin{tabular}{@{}lrrl@{}}
\toprule
Model & Parameters & Compression & One-line idea \\
\midrule
DLoc baseline (teacher) & 16.5\,M & $1\times$   & ResNet encoder $+$ 2 decoders \\
MobileNetV2 student     & 2.30\,M & $7.2\times$ & MobileNetV2 encoder $+$ upsample decoder \\
TinyCNN student         & 47\,K   & $351\times$ & 3 tiny conv blocks $+$ transpose-conv decoder \\
MobileNetV2-UNet        & 2.41\,M & $6.8\times$ & MobileNetV2 encoder $+$ U-Net skip connections \\
Mamba variant           & 653\,K  & $25\times$  & CNN stem $+$ 4 state-space blocks $+$ CNN decoder \\
\bottomrule
\end{tabular}
\end{table}

%% ==================================================================
\section{Baseline (teacher) --- recap}

The teacher is the original DLoc \code{Enc\_2Dec\_Network}: a 6-block ResNet encoder, a 9-block
location decoder, and a 12-block consistency decoder, all with InstanceNorm
($\approx16.5$\,M parameters). It is described fully in \texttt{01\_dloc\_architecture}. In the
distillation experiments it acts as the \emph{teacher}; its location-decoder output is the soft
target for the students.

%% ==================================================================
\section{MobileNetV2 student}

\textbf{Idea.} Replace the heavy ResNet encoder with a pretrained \textbf{MobileNetV2} backbone
(depthwise-separable convolutions $+$ inverted residual blocks), and decode with a lightweight
upsample-and-convolve head. The first convolution is widened from 3 (RGB) to 4 (AP) input
channels; ImageNet-pretrained weights initialize the rest.

\begin{table}[h]
\centering
\caption{MobileNetV2 student (\code{student\_mobilenet.py}). Encoder $=$ \code{mobilenet.features[:-1]}.}
\begin{tabular}{@{}lll@{}}
\toprule
Stage & Operation & Output \\
\midrule
Input        & 4-channel heatmaps                                   & $[4,161,361]$ \\
Encoder      & MobileNetV2 features (first conv $4\!\to\!32$)       & $[320, \sim\!6, \sim\!12]$ \\
Decoder 1    & Upsample$\times2$, Conv$(320\!\to\!128)$, BN, ReLU   & \\
Decoder 2    & Upsample$\times2$, Conv$(128\!\to\!64)$, BN, ReLU    & \\
Decoder 3    & Upsample$\times2$, Conv$(64\!\to\!32)$, BN, ReLU     & \\
Head         & Conv$1\times1\,(32\!\to\!1)$, Sigmoid                & \\
Resize       & Bilinear interpolate to $(161,361)$                 & $[1,161,361]$ \\
\bottomrule
\end{tabular}
\end{table}

\textbf{Rationale \& result.} Depthwise-separable convolutions factor a standard convolution into
a per-channel spatial filter plus a $1\times1$ mixing convolution, cutting compute by roughly the
kernel area. The model has $2.3$\,M parameters yet matches the teacher's in-distribution accuracy
($0.707$\,m) --- empirical evidence that DLoc is over-parameterized for this task. \textbf{Limitation:}
all spatial detail must pass through the low-resolution bottleneck, which the U-Net variant fixes.

%% ==================================================================
\section{TinyCNN student}

\textbf{Idea.} An intentionally minimal network to probe the \emph{capacity floor} of the task.
It uses one standard input convolution and then depthwise-separable blocks, downsampling twice
and upsampling twice.

\begin{table}[h]
\centering
\caption{TinyCNN student (\code{student\_tinycnn.py}). DWS $=$ depthwise-separable conv.}
\begin{tabular}{@{}lll@{}}
\toprule
Stage & Operation & Output \\
\midrule
Block 1 & Conv$(4\!\to\!16)$, BN, ReLU; DWS$(16\!\to\!16)$; MaxPool$\times2$ & $[16, H/2, W/2]$ \\
Block 2 & DWS$(16\!\to\!32)$; DWS$(32\!\to\!32)$; MaxPool$\times2$           & $[32, H/4, W/4]$ \\
Block 3 & DWS$(32\!\to\!64)$                                                 & $[64, H/4, W/4]$ \\
Dec 1   & ConvTranspose$(64\!\to\!32,\,s{=}2)$, BN, ReLU                     & $[32, H/2, W/2]$ \\
Dec 2   & ConvTranspose$(32\!\to\!16,\,s{=}2)$, BN, ReLU                     & $[16, H, W]$ \\
Head    & Conv$1\times1\,(16\!\to\!1)$, Sigmoid; resize to $(161,361)$       & $[1,161,361]$ \\
\bottomrule
\end{tabular}
\end{table}

\textbf{Result.} With only $47$\,K parameters TinyCNN reaches just $3.9$\,m median error --- it
\emph{cannot} learn both per-AP spectral feature extraction and cross-AP triangulation. It
establishes a capacity floor of $\sim$2\,M parameters and, importantly, is the model on which
\emph{distillation backfires} (Document~4).

%% ==================================================================
\section{MobileNetV2-UNet (the headline model)}

\textbf{Idea.} Keep the efficient MobileNetV2 encoder but add \textbf{U-Net skip connections}: the
encoder is split at four resolution stages and each is concatenated into the matching decoder
stage. This bypasses the information bottleneck, letting fine spatial detail from early layers
reach the output directly.

\begin{table}[h]
\centering
\caption{MobileNetV2-UNet (\code{student\_unet.py}), $2.41$\,M parameters.}
\begin{tabular}{@{}lll@{}}
\toprule
Stage & Operation & Skip out \\
\midrule
\code{inc}   & Conv$(4\!\to\!32,\,s{=}2)$, BN, ReLU6                              & --- \\
\code{enc1}  & MobileNetV2 \code{features[1:4]}                                   & 24\,ch (skip 1) \\
\code{enc2}  & MobileNetV2 \code{features[4:7]}                                   & 32\,ch (skip 2) \\
\code{enc3}  & MobileNetV2 \code{features[7:14]}                                  & 96\,ch (skip 3) \\
\code{enc4}  & MobileNetV2 \code{features[14:-1]} (bottleneck)                    & 320\,ch \\
\code{dec1}  & Up$\times2$; concat skip 3; Conv$(320{+}96\!\to\!128)$, BN, ReLU   & \\
\code{dec2}  & Up$\times2$; concat skip 2; Conv$(128{+}32\!\to\!64)$, BN, ReLU    & \\
\code{dec3}  & Up$\times2$; concat skip 1; Conv$(64{+}24\!\to\!32)$, BN, ReLU     & \\
Head         & Conv$1\times1\,(32\!\to\!1)$, Sigmoid; resize to $(161,361)$       & \\
\bottomrule
\end{tabular}
\end{table}

\textbf{Why it wins.} In a plain encoder--decoder, all spatial information is squeezed through a
low-dimensional bottleneck that discards the high-frequency detail needed for sub-metre peaks.
The skip connections re-inject encoder feature maps at three resolutions, so the decoder can
reconstruct a \emph{sharper} Gaussian peak. This yields the best in-distribution accuracy of the
family ($0.640$\,m, a $9.5\%$ improvement over the plain MobileNetV2 student) and, as a bonus,
the best noise robustness (the skip paths route around noise-corrupted bottleneck features).

%% ==================================================================
\section{Mamba variant (exploratory)}

\textbf{Idea.} Replace the ResNet bottleneck with \textbf{Mamba} selective state-space (S6) blocks
for efficient long-range sequence modelling: a CNN stem extracts local features, the 2D feature
map is flattened to a token sequence, four Mamba blocks model global relationships, and a CNN
decoder reconstructs the heatmap.

\begin{table}[h]
\centering
\caption{Mamba-DLoc (\code{mamba\_model.py}), $653$\,K parameters.}
\begin{tabular}{@{}lll@{}}
\toprule
Stage & Operation & Output \\
\midrule
Stem    & Conv$(4\!\to\!32)$; Conv$(32\!\to\!64,s{=}2)$; Conv$(64\!\to\!128,s{=}2)$ & $[128,41,91]$ \\
Tokenize& flatten spatial $\to$ sequence                                            & $[3731,128]$ \\
Core    & $4\times$ \{LayerNorm $\to$ Mamba(S6) $\to$ residual\}                     & $[3731,128]$ \\
Reshape & sequence $\to$ spatial                                                     & $[128,41,91]$ \\
Decoder & interp $\to$ Conv$(128\!\to\!64)$; interp $\to$ Conv$(64\!\to\!32)$        & $[32,161,361]$ \\
Head    & Conv$1\times1\,(32\!\to\!1)$, Sigmoid                                      & $[1,161,361]$ \\
\bottomrule
\end{tabular}
\end{table}

Each Mamba block is a selective scan: input projection, a causal depthwise Conv1d, input-dependent
SSM parameters $(\Delta, B, C)$, a discretized linear recurrence
$h_t = \bar{A}h_{t-1} + \bar{B}x_t,\; y_t = C h_t$, a skip term $D$, and a SiLU gate.
\textbf{Outcome.} The variant \emph{diverged} ($\sim$10.6\,m median): flattening a 2D heatmap into
a 1D sequence destroys spatial locality, and the model is unstable without LR warm-up and gradient
clipping. It is documented for completeness and as future work.

%% ==================================================================
\section{Knowledge distillation setup}

For the students, distillation transfers the teacher's ``soft'' spatial knowledge. To avoid
running the heavy teacher during student training, teacher outputs are \textbf{pre-computed once}
(\code{precompute\_teacher.py}, $\sim$5\,GB cached) and reused. The student loss mixes a
ground-truth term and a teacher-matching term:
\begin{equation}
\mathcal{L} \;=\; \alpha\,\text{MSE}(\hat{y},\,y)\;+\;\beta\,\text{MSE}(\hat{y},\,y_T),
\qquad \alpha = 0.7,\;\; \beta = 0.3,
\end{equation}
where $\hat y$ is the student output, $y$ the ground-truth heatmap, and $y_T$ the teacher's
output. The teacher's heatmap is \emph{broader} than the delta-like ground truth, so the second
term acts as \emph{spatial label smoothing}. Document~4 shows this helps robustness (not
in-distribution accuracy) and, below a capacity threshold, can hurt.

\paragraph{Training recipes (summary).} Students: Adam, lr $10^{-4}$, batch 64, 50 epochs.
U-Net: 120 epochs with cosine annealing. Baseline: its native schedule (StepLR, $\gamma=0.9$).
Best checkpoint is selected by minimum test-median error.

\end{document}
